% Bagean: sajarah
% Bab: biografi
% Pérangan: alonzo-church

\documentclass[../../../include/open-logic-section]{subfiles}

\begin{document}

\olfileid{his}{bio}{chu}

\olsection{Alonzo Church}

\olphoto{church-alonzo}{Alonzo Church}

Alonzo Church lair ing Washington, DC tanggal 14 Juni 1903.
Nalika isih bocah cilik, sawijining kedadéan kang nglibataké
bedhil angin njalari mripaté siji wuta. Dhèwèké ngrampungaké
sekolah persiapan ing Connecticut ing taun 1920 lan miwiti
pendhidhikan universitas ing Princeton ing taun kang padha.
Dhèwèké ngrampungaké studi doktoral ing taun 1927. Sawisé
pirang-pirang taun ana ing manca nagara, Church bali menyang
Princeton. Church kawentar banget sopan lan tliti.
Tulisané ing papan tulis rapi tanpa cacad, lan dhèwèké
ngreksa makalah-makalah wigati kanthi nutupi kanthi ati-ati
nganggo semen Duco (lem bening). Ing sanjabane kagiyatan
akadhemik, dhèwèké seneng maca majalah fiksi ilmiah lan
ora sungkan nulis marang rédaksiné yèn nemokaké
kaluputan ing tulisané.

Prestasi akadhemik Church gedhé. Bareng karo para muridé,
Stephen Kleene lan Barkley Rosser, dhèwèké ngembangaké
téori ngenani kabisan diitung kanthi efektif, yaiku
kalkulus lambda, kanthi mandiri saka pangembangan mesin
Turing déning Alan Turing. Rong dhéfinisi ngenani
kabisan diitung iku padha tegesé lan nuwuhaké apa kang
saiki diarani \emph{Tesis Church--Turing}: fungsi ing
bilangan asli bisa diitung kanthi efektif yèn lan
mung yèn bisa diitung nganggo mesin Turing (utawa
kalkulus lambda). Dhèwèké uga mbuktèkaké apa kang
saiki diarani \emph{Teorema Church}: masalah putusan
kanggo kasahihan formula orde kapisan ora bisa
dirampungaké.

Church nerusaké pakaryané nganti sepuh. Ing taun 1967
dhèwèké ninggalaké Princeton menyang UCLA, lan dadi
profèsor ing kono nganti pensiun ing taun 1990. Church
séda tanggal 1 Agustus 1995 nalika umur 92 taun.

\begin{reading}
Kanggo biografi Church kang ringkes, delengen
\citet{EndertonND}. Tulisan asli Church ngenani
kalkulus lambda lan Entscheidungsproblem
(Tesis Church) yaiku \citet{Church1936,Church1936a}.
\citet{Aspray1984} nyathet sawijining wawancara
karo Church ngenani komunitas matématika Princeton
ing dasawarsa 1930-an.
Church nulis rerangkèn resènsi buku kanggo \emph{Journal of
Symbolic Logic} wiwit taun 1936 nganti 1979. Kabèh
resènsi mau diarsipaké ing situs wèb John
MacFarlane \citep{MacFarlane2015}.
\end{reading}
\end{document}
